%%%%%%%%%%%%%%%%%%%%%%%%%%%%%%%%%%%%%%%%%%%%%%%%%%%%%%%%%%%%%%%%%%%%%%%%%%%%%%%
% APPENDIX AB: Matching and Repugnant Markets
% Market Design Without Prices
%%%%%%%%%%%%%%%%%%%%%%%%%%%%%%%%%%%%%%%%%%%%%%%%%%%%%%%%%%%%%%%%%%%%%%%%%%%%%%%
\section{Matching and Repugnant Markets: Design Without Prices}
\label{app:matching-repugnant}

%==============================================================================
% PART 0: INTEGRATION OVERVIEW
%==============================================================================
\subsection{Framework Integration Map}

This appendix demonstrates how matching theory and repugnant market design 
provide essential foundations for the $(C, \Psi, C^*, K, Q)$ framework in 
contexts where price mechanisms cannot or should not operate. Many of the 
most important allocation problems---organ transplants, school assignments, 
medical residencies---require reaching $C^*$ without monetary transfers. 
The work of Alvin Roth, Lloyd Shapley, and their collaborators shows how 
to design $\Psi_{mechanism}$ to achieve stable, efficient, and fair allocations.

\begin{center}
\renewcommand{\arraystretch}{1.4}
\begin{tabular}{|l|l|l|}
\hline
\textbf{Framework Element} & \textbf{Matching Theory Contribution} & \textbf{Key Papers} \\
\hline
$C^*_{stable}$ (No blocking pairs) & Stable matching theory & Gale-Shapley 1962 \\
$\Psi_{mechanism}$ design & Deferred acceptance & Roth 1984, Roth-Peranson 1999 \\
$\Psi_{moral}$ constraints & Repugnance theory & Roth 2007 \\
$C^*$ without transfers & Kidney exchange & Roth-Sönmez-Ünver 2004 \\
$\Psi_{mechanism}$ comparison & School choice & Abdulkadiroğlu-Sönmez 2003 \\
$C^*$ timing & Unraveling & Niederle-Roth 2003 \\
\hline
\end{tabular}
\end{center}

\paragraph{Matching Theory's Unique Role.}
Matching theory contributes uniquely to the framework:
\begin{enumerate}
    \item \textbf{Non-price allocation}: How to reach $C^*$ without money
    \item \textbf{Stability concept}: $C^*_{stable}$ as no blocking pairs
    \item \textbf{Mechanism design}: $\Psi_{mechanism}$ that induces truth-telling
    \item \textbf{Repugnance}: $\Psi_{moral}$ as binding constraint
    \item \textbf{Practical implementation}: Theory that saves lives
\end{enumerate}

%==============================================================================
% PART I: STABLE MATCHING THEORY
%==============================================================================
\subsection{Part I: Stable Matching Theory – The Foundation}

\subsubsection{The Marriage Problem}

Gale \& Shapley (1962) formalized matching:

\begin{equation}
\boxed{
\text{Given preferences of men over women and women over men, find a stable matching}
}
\label{eq:gs-problem}
\end{equation}

\paragraph{Stability Definition.}
A matching $\mu$ is \textbf{stable} if there is no blocking pair---no man $m$ 
and woman $w$ who both prefer each other to their current partners:

\begin{equation}
\nexists (m, w): w \succ_m \mu(m) \text{ and } m \succ_w \mu(w)
\end{equation}

\subsubsection{Framework Translation}

Stable matching \textit{is} $C^*$ for two-sided markets:

\begin{equation}
\boxed{
C^*_{stable} = \{\mu : \text{No blocking pairs}\}
}
\label{eq:stable-cstar}
\end{equation}

\paragraph{Key Insight.}
Stability is self-enforcing equilibrium: no pair has incentive to deviate.
This is $C^*$ in the matching game.

\subsubsection{The Deferred Acceptance Algorithm}

Gale-Shapley proved existence constructively:

\begin{algorithm}[Deferred Acceptance (Men-Proposing)]
\begin{enumerate}
    \item Each man proposes to his top choice
    \item Each woman tentatively accepts best proposal, rejects others
    \item Rejected men propose to next choice
    \item Repeat until no rejections
\end{enumerate}
\end{algorithm}

\paragraph{Key Results.}
\begin{itemize}
    \item Algorithm terminates in finite steps
    \item Output is always stable
    \item Men-proposing DA is optimal for men (worst stable for women)
    \item Women-proposing DA is optimal for women
\end{itemize}

\subsubsection{Framework Translation}

DA is a $\Psi_{mechanism}$ that computes $C^*$:

\begin{equation}
\Psi_{DA} \to C^*_{stable} \quad \text{(guaranteed)}
\end{equation}

The mechanism \textit{implements} the stable matching---a remarkable 
achievement where theory directly produces practical solutions.

\subsubsection{The Lattice Structure}

The set of stable matchings forms a lattice:

\begin{equation}
C^*_{men-optimal} \succeq_{men} C^* \succeq_{men} C^*_{women-optimal}
\end{equation}

All agents on one side agree on which stable matching is best for their side.

%==============================================================================
% PART II: MARKET DESIGN PRINCIPLES
%==============================================================================
\subsection{Part II: Market Design Principles}

\subsubsection{Roth's Three Conditions}

Alvin Roth identified three conditions for well-functioning markets:

\begin{equation}
\boxed{
\text{Functioning Market} \Leftrightarrow \text{Thickness} + \text{No Congestion} + \text{Safety}
}
\label{eq:roth-three}
\end{equation}

\paragraph{1. Thickness.}
Enough participants at the same time:
\begin{equation}
\Psi_{thickness} = |\text{Participants}| \times \Pr(\text{Simultaneous})
\end{equation}

Thin markets lack good matches; thick markets enable efficient $C^*$.

\paragraph{2. No Congestion.}
Enough time to evaluate options:
\begin{equation}
\Psi_{congestion}^{low} \Rightarrow \text{Time to process offers, compare options}
\end{equation}

Congested markets force hasty, suboptimal decisions.

\paragraph{3. Safety.}
Incentives for honest participation:
\begin{equation}
\Psi_{safe} \Rightarrow \text{Truth-telling is optimal strategy}
\end{equation}

Unsafe markets encourage strategic manipulation, undermining efficiency.

\subsubsection{Framework Translation}

Market design is choosing $\Psi_{mechanism}$ to satisfy all three:
\begin{equation}
\Psi^*_{mechanism} = \arg\max_{\Psi} Q(\Psi) \quad \text{s.t. Thick, Uncongested, Safe}
\end{equation}

\subsubsection{Strategy-Proofness}

A mechanism is \textbf{strategy-proof} if truth-telling is dominant strategy:

\begin{equation}
\boxed{
u_i(\text{True preferences}) \geq u_i(\text{Any manipulation}) \quad \forall i
}
\label{eq:strategyproof}
\end{equation}

\paragraph{Key Result.}
Deferred acceptance is strategy-proof for the proposing side:
\begin{itemize}
    \item Men-proposing DA: Truth-telling optimal for men
    \item Women-proposing DA: Truth-telling optimal for women
\end{itemize}

\paragraph{Framework Translation.}
Strategy-proofness means $\Psi_{mechanism}$ aligns individual and social $C^*$:
\begin{equation}
C^*_{individual} = C^*_{social} \quad \text{(no conflict)}
\end{equation}

%==============================================================================
% PART III: REPUGNANT MARKETS
%==============================================================================
\subsection{Part III: Repugnant Markets – When Prices Cannot Work}

\subsubsection{The Concept of Repugnance}

Roth (2007) identifies \textbf{repugnance} as market constraint:

\begin{equation}
\boxed{
\Psi_{moral} \Rightarrow \text{Price mechanism prohibited}
}
\label{eq:repugnance}
\end{equation}

\paragraph{Definition.}
Some transactions that would be mutually beneficial are blocked because 
third parties find them repugnant---morally objectionable.

\subsubsection{Examples of Repugnant Markets}

\begin{center}
\renewcommand{\arraystretch}{1.3}
\begin{tabular}{|l|l|l|}
\hline
\textbf{Market} & \textbf{Repugnance} & \textbf{Status} \\
\hline
Organ sales & Commodification of body & Prohibited almost everywhere \\
Prostitution & Commodification of intimacy & Prohibited in most places \\
Surrogacy & Commodification of reproduction & Varies by jurisdiction \\
Dwarf tossing & Dignity violation & Prohibited in France \\
Horse meat & Cultural taboo & Prohibited in US, legal in EU \\
Egg/sperm donation & Partial commodification & Compensation allowed \\
Blood donation & Varies by culture & Paid in US, unpaid in UK \\
\hline
\end{tabular}
\end{center}

\subsubsection{Framework Translation}

Repugnance is $\Psi_{moral}$ constraining mechanism choice:
\begin{equation}
\Psi_{mechanism} \in \{\text{Non-price mechanisms}\} \quad \text{when } \Psi_{moral} \text{ binds}
\end{equation}

\paragraph{Key Insight.}
The framework must accommodate non-price allocation:
\begin{equation}
C^*_{repugnant} = \arg\max Q \quad \text{s.t. } \Psi_{moral}, \text{ no prices}
\end{equation}

This is why matching theory matters: it provides tools for $C^*$ without money.

\subsubsection{Repugnance is Context-Dependent}

What is repugnant varies across:
\begin{itemize}
    \item \textbf{Time}: Life insurance was once considered repugnant
    \item \textbf{Space}: Horse meat legal in France, not US
    \item \textbf{Culture}: Attitudes toward surrogacy vary widely
\end{itemize}

\paragraph{Framework Translation.}
$\Psi_{moral}$ is itself context-dependent:
\begin{equation}
\Psi_{moral} = f(\Psi_{culture}, \Psi_{time}, \Psi_{place})
\end{equation}

Market designers must work within these constraints, not against them.

%==============================================================================
% PART IV: KIDNEY EXCHANGE
%==============================================================================
\subsection{Part IV: Kidney Exchange – $C^*$ Without Money}

\subsubsection{The Problem}

Kidney transplants save lives, but:
\begin{itemize}
    \item Selling kidneys is illegal (repugnance)
    \item Many willing donors are incompatible with intended recipient
    \item Waiting list has 100,000+ people; ~5,000 die annually
\end{itemize}

\subsubsection{The Solution: Exchange}

Roth, Sönmez \& Ünver (2004, 2007) designed kidney exchange:

\begin{equation}
\boxed{
\text{Pair } (D_1 \to R_2) + \text{Pair } (D_2 \to R_1) = \text{Two transplants}
}
\label{eq:kidney-exchange}
\end{equation}

\paragraph{The Mechanism.}
\begin{itemize}
    \item Patient-donor pairs register incompatibility
    \item Algorithm finds cycles and chains of compatible exchanges
    \item Simultaneous surgery prevents reneging
\end{itemize}

\subsubsection{Framework Translation}

Kidney exchange is $C^*$ without transfers:
\begin{equation}
C^*_{kidney} = \arg\max \text{Transplants} \quad \text{s.t. } \Psi_{moral}, \text{ no payments}
\end{equation}

\paragraph{Key Innovation.}
In-kind exchange avoids repugnance while achieving efficiency:
\begin{equation}
\Psi_{exchange} \not\in \Psi_{repugnant} \quad \text{(exchange is acceptable)}
\end{equation}

\subsubsection{Chains and Non-Directed Donors}

Extension to chains initiated by altruistic donors:

\begin{equation}
D_0 \to R_1, \quad D_1 \to R_2, \quad D_2 \to R_3, \quad \ldots
\end{equation}

\paragraph{Non-Simultaneous Extended Altruistic Donor (NEAD) Chains.}
\begin{itemize}
    \item Altruistic donor starts chain
    \item Chain can extend over time
    \item ``Bridge donor'' waits for next compatible recipient
\end{itemize}

\subsubsection{Impact}

Kidney exchange has:
\begin{itemize}
    \item Facilitated thousands of transplants
    \item Expanded from pairs to chains to multi-hospital networks
    \item Been adopted in US, UK, Netherlands, and beyond
\end{itemize}

\paragraph{Framework Translation.}
Theory directly increases $Q$:
\begin{equation}
Q_{kidney exchange} = \text{Lives saved} \times \text{Life-years gained}
\end{equation}

%==============================================================================
% PART V: SCHOOL CHOICE
%==============================================================================
\subsection{Part V: School Choice – Mechanism Comparison}

\subsubsection{The Problem}

Assigning students to public schools:
\begin{itemize}
    \item Students have preferences over schools
    \item Schools have priorities (siblings, proximity, lottery)
    \item No prices; assignment by mechanism
\end{itemize}

\subsubsection{The Boston Mechanism}

Traditional Boston mechanism:
\begin{enumerate}
    \item Students rank schools
    \item Round 1: Assign to first choice if capacity allows
    \item Round 2: Remaining students assigned to second choice
    \item Continue until all assigned
\end{enumerate}

\paragraph{Problem.}
Boston mechanism is \textbf{not strategy-proof}:
\begin{equation}
\text{Optimal strategy} \neq \text{True preferences}
\end{equation}

Students should avoid ranking popular schools they're unlikely to get.
This disadvantages unsophisticated families.

\subsubsection{Deferred Acceptance for Schools}

Abdulkadiroğlu \& Sönmez (2003) proposed student-proposing DA:

\begin{equation}
\boxed{
\Psi_{DA} \text{ is strategy-proof for students}
}
\label{eq:school-da}
\end{equation}

\paragraph{Advantages.}
\begin{itemize}
    \item Truth-telling is optimal for all students
    \item No gaming required; levels playing field
    \item Outcome is stable (no justified envy)
\end{itemize}

\subsubsection{New York and Boston Reforms}

Abdulkadiroğlu, Pathak \& Roth redesigned:
\begin{itemize}
    \item \textbf{NYC (2003)}: Replaced chaotic system with DA for 90,000+ students
    \item \textbf{Boston (2005)}: Replaced Boston mechanism with DA
\end{itemize}

\paragraph{Results.}
\begin{itemize}
    \item More students get top choices
    \item Less strategic behavior observed
    \item More equitable across socioeconomic groups
\end{itemize}

\subsubsection{Framework Translation}

School choice shows mechanism matters for $C^*$ quality:
\begin{equation}
Q(\Psi_{DA}) > Q(\Psi_{Boston}) \quad \text{(welfare improvement)}
\end{equation}

Same preferences, same priorities, different mechanism $\Rightarrow$ different $C^*$.

%==============================================================================
% PART VI: UNRAVELING AND TIMING
%==============================================================================
\subsection{Part VI: Unraveling and Market Timing}

\subsubsection{The Unraveling Problem}

Many markets suffer from \textbf{unraveling}---transactions occurring 
progressively earlier:

\begin{equation}
\boxed{
t_{transaction} \to 0 \quad \text{(timing moves earlier and earlier)}
}
\label{eq:unraveling}
\end{equation}

\paragraph{Examples.}
\begin{itemize}
    \item Law clerk hiring: Offers during 2L year, then 1L, then before law school
    \item College football bowls: Bids in October for January games
    \item Sorority rush: Compressed to single weekend
\end{itemize}

\subsubsection{Why Unraveling Happens}

Strategic incentives drive early offers:
\begin{itemize}
    \item \textbf{Employers}: Lock in good candidates before competitors
    \item \textbf{Candidates}: Accept early offer to avoid risk of nothing
\end{itemize}

This creates prisoner's dilemma: all worse off but no one can unilaterally stop.

\subsubsection{Framework Translation}

Unraveling is $C^*$ timing failure:
\begin{equation}
C^*_{early} \prec C^*_{late} \quad \text{but} \quad C^*_{early} \text{ is equilibrium}
\end{equation}

Inefficient early matching is stable because of coordination failure.

\subsubsection{Niederle \& Roth: Gastroenterology}

Niederle \& Roth (2003) studied gastroenterology fellowship market:
\begin{itemize}
    \item Market had unraveled to 2 years before position start
    \item Match introduced, then collapsed, then reinstated
    \item Shows fragility of coordinated timing
\end{itemize}

\paragraph{Framework Translation.}
Market thickness requires coordination on timing:
\begin{equation}
\Psi_{thickness} = f(\text{Coordinated timing})
\end{equation}

Without coordination mechanism, markets thin and unravel.

\subsubsection{Solutions to Unraveling}

\begin{center}
\renewcommand{\arraystretch}{1.3}
\begin{tabular}{|l|l|}
\hline
\textbf{Solution} & \textbf{Mechanism} \\
\hline
Centralized matching & NRMP for medical residencies \\
Binding agreements & NCAA signing dates \\
Information control & Delay release of information \\
Exploding offers ban & Prohibit short deadlines \\
\hline
\end{tabular}
\end{center}

%==============================================================================
% PART VII: SYNTHESIS
%==============================================================================
\subsection{Part VII: Synthesis – Matching and Market Design Equations}

\subsubsection{Complete Framework Translation}

\begin{center}
\renewcommand{\arraystretch}{1.5}
\begin{tabular}{|l|c|l|}
\hline
\textbf{Finding} & \textbf{Equation} & \textbf{Framework Element} \\
\hline
Stable matching & $\nexists$ blocking pair & $C^*_{stable}$ definition \\
\hline
DA algorithm & Computes $C^*$ & $\Psi_{mechanism}$ implementation \\
\hline
Strategy-proofness & Truth = optimal & $C^*_{individual} = C^*_{social}$ \\
\hline
Repugnance & $\Psi_{moral}$ binds & Prices prohibited \\
\hline
Kidney exchange & In-kind $C^*$ & $C^*$ without transfers \\
\hline
School choice & $Q(\Psi_{DA}) > Q(\Psi_{Boston})$ & Mechanism matters \\
\hline
Unraveling & $t \to 0$ & $C^*$ timing failure \\
\hline
Three conditions & Thick, uncongested, safe & $\Psi$ for functioning markets \\
\hline
\end{tabular}
\end{center}

\subsubsection{Matching Theory's Unique Contribution}

\textbf{Without Matching Theory, the framework would:}
\begin{itemize}
    \item Assume prices always available
    \item Miss stability as equilibrium concept
    \item Lack tools for non-price allocation
    \item Not understand repugnance constraints
    \item Miss timing/thickness problems
\end{itemize}

\textbf{With Matching Theory, the framework gains:}
\begin{itemize}
    \item Non-price $C^*$: Matching without money
    \item Stability: Self-enforcing allocations
    \item Mechanism design: $\Psi$ that implements $C^*$
    \item Repugnance: $\Psi_{moral}$ as constraint
    \item Practical impact: Lives saved, welfare improved
\end{itemize}

%==============================================================================
% PART VIII: COMPLETE PAPER DOCUMENTATION
%==============================================================================
\subsection{Part VIII: Complete Paper Documentation}

%--- FOUNDATIONAL THEORY ---
\subsubsection{Foundational Theory}

\paragraph{Paper 1: Gale \& Shapley (1962)}

``College Admissions and the Stability of Marriage.''
\textit{American Mathematical Monthly}, 69(1), 9--15.

\textbf{Summary.}
Introduces stable matching problem and deferred acceptance algorithm.
Proves existence of stable matching for any preferences.

\textbf{Framework Translation.}
\begin{equation}
C^*_{stable} \text{ exists and DA computes it}
\end{equation}

\textbf{Key Finding.}
Nobel Prize (Shapley 2012). Foundation of matching theory. 10,000+ citations.

\paragraph{Paper 2: Shapley \& Scarf (1974)}

``On Cores and Indivisibility.''
\textit{Journal of Mathematical Economics}, 1(1), 23--37.

\textbf{Summary.}
Housing market model with indivisible goods. Shows core exists and equals 
competitive allocation. Top Trading Cycles algorithm.

\textbf{Framework Translation.}
\begin{equation}
C^*_{core} = C^*_{competitive} \quad \text{for housing markets}
\end{equation}

\paragraph{Paper 3: Roth (1984)}

``The Evolution of the Labor Market for Medical Interns and Residents: 
A Case Study in Game Theory.''
\textit{Journal of Political Economy}, 92(6), 991--1016.

\textbf{Summary.}
Shows NRMP (National Resident Matching Program) is equivalent to hospital-proposing 
DA. Explains historical evolution through game-theoretic lens.

\textbf{Framework Translation.}
$\Psi_{NRMP}$ implements $C^*_{stable}$ for medical labor market.

\textbf{Key Finding.}
Nobel Prize (Roth 2012). Connected theory to practice.

%--- MARKET DESIGN ---
\subsubsection{Market Design Practice}

\paragraph{Paper 4: Roth \& Peranson (1999)}

``The Redesign of the Matching Market for American Physicians: 
Some Engineering Aspects of Economic Design.''
\textit{American Economic Review}, 89(4), 748--780.

\textbf{Summary.}
Describes redesign of NRMP to handle couples and other complexities.
Shows how theory guides practical mechanism design.

\textbf{Framework Translation.}
$\Psi_{mechanism}$ design for real-world constraints.

\paragraph{Paper 5: Roth (2008)}

``What Have We Learned from Market Design?''
\textit{Economic Journal}, 118(527), 285--310.

\textbf{Summary.}
Synthesizes lessons from market design practice. Thickness, congestion, 
safety as key conditions.

\textbf{Framework Translation.}
Three conditions for $\Psi_{mechanism}$ to achieve efficient $C^*$.

\paragraph{Paper 6: Roth (2015)}

\textit{Who Gets What---and Why: The New Economics of Matchmaking and Market Design.}
Houghton Mifflin Harcourt.

\textbf{Summary.}
Accessible book on market design. Covers kidney exchange, school choice, 
job markets, dating.

\textbf{Framework Translation.}
Complete accessible treatment of matching and $\Psi_{mechanism}$ design.

\paragraph{Paper 7: Roth (2018)}

``Marketplaces, Markets, and Market Design.''
\textit{American Economic Review}, 108(7), 1609--1658.

\textbf{Summary.}
Presidential address synthesizing market design. Distinguishes marketplaces 
(venues) from markets (price-mediated exchange).

\textbf{Framework Translation.}
\begin{equation}
\Psi_{marketplace} \neq \Psi_{market} \quad \text{(different design problems)}
\end{equation}

%--- REPUGNANCE ---
\subsubsection{Repugnance}

\paragraph{Paper 8: Roth (2007)}

``Repugnance as a Constraint on Markets.''
\textit{Journal of Economic Perspectives}, 21(3), 37--58.

\textbf{Summary.}
Introduces repugnance as economic concept. Documents variation across 
transactions, time, and place.

\textbf{Framework Translation.}
\begin{equation}
\Psi_{moral} \Rightarrow \text{Price mechanisms prohibited}
\end{equation}

\textbf{Key Finding.}
Foundational paper on repugnance. 1,500+ citations.

%--- KIDNEY EXCHANGE ---
\subsubsection{Kidney Exchange}

\paragraph{Paper 9: Roth, Sönmez \& Ünver (2004)}

``Kidney Exchange.''
\textit{Quarterly Journal of Economics}, 119(2), 457--488.

\textbf{Summary.}
Introduces kidney exchange mechanism. Shows how to organize exchanges 
among incompatible patient-donor pairs.

\textbf{Framework Translation.}
\begin{equation}
C^*_{kidney} = \text{Maximum transplants without money}
\end{equation}

\textbf{Key Finding.}
Launched kidney exchange programs. 2,000+ citations.

\paragraph{Paper 10: Roth, Sönmez \& Ünver (2007)}

``Efficient Kidney Exchange: Coincidence of Wants in Markets with 
Compatibility-Based Preferences.''
\textit{American Economic Review}, 97(3), 828--851.

\textbf{Summary.}
Extends to larger cycles and chains. Shows efficient mechanisms for 
different constraints.

\textbf{Framework Translation.}
Optimal $\Psi_{mechanism}$ depends on operational constraints.

\paragraph{Paper 11: Ashlagi \& Roth (2014)}

``Free Riding and Participation in Large Scale, Multi-Hospital Kidney Exchange.''
\textit{Theoretical Economics}, 9(3), 817--863.

\textbf{Summary.}
Studies incentives for hospitals to participate fully in exchange.
Shows potential for free-riding in multi-hospital systems.

\textbf{Framework Translation.}
\begin{equation}
C^*_{participation} \neq C^*_{efficient} \quad \text{(participation incentives)}
\end{equation}

%--- SCHOOL CHOICE ---
\subsubsection{School Choice}

\paragraph{Paper 12: Abdulkadiroğlu \& Sönmez (2003)}

``School Choice: A Mechanism Design Approach.''
\textit{American Economic Review}, 93(3), 729--747.

\textbf{Summary.}
Applies mechanism design to school choice. Compares Boston mechanism 
to DA. Shows DA is strategy-proof, Boston is not.

\textbf{Framework Translation.}
\begin{equation}
Q(\Psi_{DA}) > Q(\Psi_{Boston})
\end{equation}

\textbf{Key Finding.}
Foundation for school choice reform. 3,500+ citations.

\paragraph{Paper 13: Abdulkadiroğlu, Pathak \& Roth (2005)}

``The New York City High School Match.''
\textit{American Economic Review P\&P}, 95(2), 364--367.

\textbf{Summary.}
Describes design and implementation of NYC high school match for 
90,000+ students annually.

\textbf{Framework Translation.}
Large-scale $\Psi_{mechanism}$ implementation.

\paragraph{Paper 14: Abdulkadiroğlu, Pathak, Roth \& Sönmez (2005)}

``The Boston Public School Match.''
\textit{American Economic Review P\&P}, 95(2), 368--371.

\textbf{Summary.}
Describes redesign of Boston school assignment from Boston mechanism to DA.

\textbf{Framework Translation.}
Mechanism replacement improves $C^*$ quality.

\paragraph{Paper 15: Pathak \& Sönmez (2008)}

``Leveling the Playing Field: Sincere and Sophisticated Players in the 
Boston Mechanism.''
\textit{American Economic Review}, 98(4), 1636--1652.

\textbf{Summary.}
Shows Boston mechanism hurts unsophisticated families who report true 
preferences against sophisticated families who game.

\textbf{Framework Translation.}
\begin{equation}
\Psi_{Boston} \text{ favors sophisticated} \Rightarrow \text{Inequitable } C^*
\end{equation}

%--- UNRAVELING ---
\subsubsection{Unraveling and Timing}

\paragraph{Paper 16: Niederle \& Roth (2003)}

``Unraveling Reduces the Scope of an Entry Level Labor Market: 
Gastroenterology with and without a Centralized Match.''
\textit{Journal of Political Economy}, 111(6), 1342--1352.

\textbf{Summary.}
Studies gastroenterology fellowship market as match collapsed and 
market unraveled. Shows unraveling reduces mobility and match quality.

\textbf{Framework Translation.}
\begin{equation}
\text{Unraveling} \Rightarrow \Psi_{thickness} \downarrow \Rightarrow Q \downarrow
\end{equation}

\paragraph{Paper 17: Niederle \& Roth (2009)}

``Market Culture: How Rules Governing Exploding Offers Affect Market 
Performance.''
\textit{American Economic Journal: Microeconomics}, 1(2), 199--219.

\textbf{Summary.}
Experimental study of exploding offers. Shows short deadlines reduce 
market efficiency and participant welfare.

\textbf{Framework Translation.}
$\Psi_{timing}$ affects $C^*$ quality.

%--- EXTENSIONS ---
\subsubsection{Extensions and Theory}

\paragraph{Paper 18: Hatfield \& Milgrom (2005)}

``Matching with Contracts.''
\textit{American Economic Review}, 95(4), 913--935.

\textbf{Summary.}
Unifies matching and auction theory. Shows many results extend to 
matching with contracts (job market with wages).

\textbf{Framework Translation.}
\begin{equation}
C^*_{matching} \text{ generalizes to contracts and prices}
\end{equation}

\textbf{Key Finding.}
Major theoretical unification. 2,000+ citations.

\paragraph{Paper 19: Sönmez \& Ünver (2010)}

``House Allocation with Existing Tenants: A Characterization.''
\textit{Games and Economic Behavior}, 69(2), 425--445.

\textbf{Summary.}
Characterizes mechanisms for housing allocation when some agents have 
existing assignments. Balances efficiency and property rights.

\textbf{Framework Translation.}
$C^*$ with endowments requires respecting status quo.

\paragraph{Paper 20: Kominers \& Sönmez (2016)}

``Matching with Slot-Specific Priorities: Theory.''
\textit{Theoretical Economics}, 11(2), 683--710.

\textbf{Summary.}
Extends matching to handle slot-specific priorities (different criteria 
for different seats). Applications to cadet-branch matching.

\textbf{Framework Translation.}
$\Psi_{mechanism}$ for complex priority structures.

%%%%%%%%%%%%%%%%%%%%%%%%%%%%%%%%%%%%%%%%%%%%%%%%%%%%%%%%%%%%%%%%%%%%%%%%%%%%%%%
\subsection{Citation and Verification Note}

\textbf{Nobel Prizes represented:}
\begin{itemize}
    \item Shapley \& Roth (2012) ``for the theory of stable allocations and 
          the practice of market design''
\end{itemize}

\textbf{Impact.}
\begin{itemize}
    \item NRMP matches 40,000+ medical residents annually
    \item NYC matches 90,000+ students annually
    \item Kidney exchange has facilitated thousands of transplants
    \item Theory directly saves lives
\end{itemize}

Combined Google Scholar citations for the 20 papers: $>$50,000.

This appendix covers matching theory from Gale-Shapley through modern 
market design, providing foundations for non-price allocation within 
the $(C, \Psi, C^*, K, Q)$ framework.

%%%%%%%%%%%%%%%%%%%%%%%%%%%%%%%%%%%%%%%%%%%%%%%%%%%%%%%%%%%%%%%%%%%%%%%%%%%%%%%
